\documentclass[../master.tex]{subfiles}

\begin{document}
\chapter{一般化パウリ演算子とqudit}
誤り訂正符号の世界はほとんどがmod 2の世界で記述されている。
しかし、mod 2の世界というのは0と1しかない世界であるため、
mod 2の特異性からくる符号の性質なのか、もっと一般的な性質なのかがわからないことがある。
この章では、mod 2を一般化してmod \(N\)の世界での量子誤り訂正符号を記述するための道具を整備し、
その道具を用いてShor符号を記述することを目標とする。

\section{mod Nの代数}
符号理論においては、群や体といった代数学の言葉を用いて符号の性質を記述することが多い。
そこで、このセクションでは代数学の入門として、mod \(N\)の数の世界を紹介する。

古典誤り訂正符号の例として0,1のビットによる繰り返し符号を考えてきた。
これらの誤り検出に使う検査行列によるシンドロームの計算を改めて見てみる。
\begin{equation}\begin{aligned}[b]
\begin{pmatrix}
    1 \\ 0
\end{pmatrix}=
\begin{pmatrix}
    1 & 1 & 0 \\
    0 & 1 & 1
\end{pmatrix}
\begin{pmatrix}
    0 \\
    1 \\
    1
\end{pmatrix}
\end{aligned}\end{equation}
n番目のビットを\(x_n\)の様に表すとする。
この検査行列として本来行いたいのは、
1番目と2番目のビットが同じかどうか、2番目と3番目のビットが同じかどうかを調べることである。
これを数式として表すと、\(x_1-x_2=0\), \(x_2-x_3=0\)となっているかを計算することになる。
しかし、実際のこの計算を見てみると、
検査行列の1行目では\(x_1+x_2\),
検査行列の2行目では\(x_2+x_3\)を計算していることがわかる。
これは何に由来しているかというと、mod 2ならではの性質である。
mod 2の世界の和は次の様に表される。
\begin{equation}\begin{aligned}[b]
    0 + 0 &= 0, & 0 + 1 &= 1, & 1 + 0 &= 1, & 1 + 1 &= 0
\end{aligned}\end{equation}
この式のうち\(1+1=0\)に注目すると、
\(-1 = 1\)となるため、mod 2の世界では和と差の区別がつかなくなってしまう。
このようなmod 2の特異性を避けるために、2ではなく一般的なNを用いてmod \(N\)の世界を考えることがある。
mod \(N\) の世界の数学的な定義のため、すこし回り道をしよう。

数学に限らず、本来は違うもの同士を一定の規則に従って同一視したいことはよくある。
そうしたもの同士が同一視される基準として、以下の3つの条件を満たしているべきだろう。
\begin{NoteBox}[title=同値関係]{}
    ある集合\(S\)の要素\(a,b,c\)に対して、
    これらが同一視できるとき\(a\sim b\)のように書く。
    \begin{itemize}
        \item 反射律: \(a \sim a\).
        \item 対称律: \(a \sim b\)ならば、\(b \sim a\).
        \item 推移律: \(a \sim b\)かつ\(b \sim c\)ならば\(a \sim c\)
    \end{itemize}
    これらの性質が成り立つ二項関係\(\sim\)は、
    数学の言葉で同値関係であると言う。
    \footnote{同値関係と等号はよく似ているがこれらは違うものである。
    等号はこれら3つの性質に加えて、代入可能であるという性質も満たす必要がある。}
\end{NoteBox}

整数を使って同値関係の例を考えてみよう。
整数の集合は\(\mathbb{Z}\), \(N\)の倍数の集合のこと\(N\mathbb{Z}\)と書く。
整数\(a,b\)に対して、二項関係\(a\sim b\)を\(a\)と\(b\)の差が\(N\)の倍数である、つまり
\begin{equation}\begin{aligned}[b]
    a \sim b \quad \xleftrightarrow{\text{def}} \quad a - b \in N\mathbb{Z}
\end{aligned}\end{equation}
と定義する。
このようにして定義した\(\sim\)は先ほどの3条件をどれもみたすことが確認できる同値関係である。

元々違うもの同士の同一視を行いたいために同値関係を導入したのであった。
ではこの整数の例によって同一視されたもの同士で集合を作ってみよう。
すると
\begin{equation}\begin{aligned}[b]
    \{0,\, N,\,-N,\, 2N, \ldots\}&=\{mN\mid m\in \mathbb{Z}\}=:[0]\\
    \{1,\, N+1,\,-N+1,\, 2N+1, \ldots\}&=\{1+mN\mid m\in \mathbb{Z}\}=:[1]\\
    &\vdots \\
    \{N-1,\, -N-1,\,2N-1, \ldots\}&=\{N-1+mN\mid m\in \mathbb{Z}\}=:[N-1]
\end{aligned}\end{equation}
のように、\(N\)で割った余りが同じ整数同士で同一視されていることがわかる。
このように同値関係によって集合の要素同士が繋がっていて、
繋がりかたで元の集合はグループ分けすることができる。
このグループ1つ1つのことを同値類と呼び、同値類全体の集合のことを商集合と呼ぶ。
同値類と商集合の定義を一般的に書いてみよう。

\begin{NoteBox}[title=同値類と商集合]{}
    ある集合\(S\)と同値関係\(\sim\)が与えられたとき、\(a\in S\)に対して、
    \begin{equation}\begin{aligned}[b]
        [a] = \{s\in S \mid a \sim s\}
    \end{aligned}\end{equation}
    と定義する。このようにして定義された\([a]\)のことを\(a\)の同値類と呼ぶ。
    また、同値類を表すために選ばれた要素\(a\)のことを同値類の代表元と呼ぶ。
    同値類全体の集合のことを商集合と呼び、
    \begin{equation}\begin{aligned}[b]
        S/\sim\, = \{[a] \mid a\in S\}
    \end{aligned}\end{equation}
    と書く。

    また、同値関係として集合\(S\)の元同士の差が
    \(S\)の部分集合\(R\)の元であることを定義することが実用上よく行われる。
    つまり、ある集合\(S\)とその部分集合\(R\)
    \footnote{正確に書くと\(S\)は加法についての可換群、\(R\)はその部分群}
    が与えられたとき、\(a,b\in S\)に対して、
    \begin{equation}\begin{aligned}[b]
        a \sim_R b \quad \xleftrightarrow{\text{def}} \quad a - b \in R
    \end{aligned}\end{equation}
    と同値関係を定義する。
    このようにして定義された同値関係によってできる商集合のことを\(S/R\)と書く。
    また、この商集合の元である同値類を
    \begin{equation}\begin{aligned}[b]
        [a] = a+R = \{a+r \mid r\in R\}
    \end{aligned}\end{equation}
    のように書くこともある。
\end{NoteBox}

整数での例では、同値類は\(N\)で割った余りが同じ整数同士である\([0], [1], \ldots, [N-1]\)それぞれのことである。
これらの代表元は\(0, 1, \ldots, N-1\)、
この商集合は\(\{[0],[1],\ldots, [N-1]\}\)のことである。
つまり、整数を\(N\)で割った余りが同じ整数同士を同一視することによって、
整数の集合\(\mathbb{Z}\)は\(N\)個の同値類に分割される。
この同値関係は\(a \sim b\)を\(a\)と\(b\)の差が\(N\)の倍数であることと定義しているので、
商集合は\(\mathbb{Z}/\sim\)ではなく\(\mathbb{Z}/N\mathbb{Z}\)と書かれる。
さらにこれは頻繁に使われることから、\(\mathbb{Z}_N := \mathbb{Z}/N\mathbb{Z}\)と書かれる。

同値類がおおよそ mod \(N\) の世界の数でありそうなことはわかった。
mod \(N\) をとった後でも、元同士の和や積を考えていたわけだが、
同値類は集合の一種であるため、同値類同士の和や積を定義する必要がある。
\(\mathbb{Z}_N\)における和と積は
\begin{equation}\begin{aligned}[b]
    [a]+[b] &:=[a+b], &
    [a][b] &:=[ab]
\end{aligned}\end{equation}
つまり
\begin{equation}\begin{aligned}[b]
    a+N\mathbb{Z} + b+N\mathbb{Z} &:= (a+b)+N\mathbb{Z}, &
    (a+N\mathbb{Z})(b+N\mathbb{Z}) &:= ab + N\mathbb{Z}
\end{aligned}\end{equation}
のように定義される。
また、差が\(N\)の倍数であれば同一視されるということは、
代表元同士の差が\(N\)の倍数であれば同一視されるということでもある。
つまり、
\begin{equation}\begin{aligned}[b]
    [N-1] = (N-1)+ N\mathbb{Z} = -1 +N\mathbb{Z} = [-1]
\end{aligned}\end{equation}
のように書いてもよく、代表元は一意に決まっていない。
この冗長性により、\((a+b)\)や\((ab)\)と書かれている代表元の中から、
0から\(N-1\)の間にあるものを選べば、整数を\(N\)で割ったあまりについての計算規則だとみなすことができる
\footnote{この都合の良い代表元を選んでもよいというのは、数学の専門用語である "well-defined" であるという意味である。
この概念はいろんな表示がありうるような対象を作った際に不整合な関係を持っていないことを示す。
これを確認するのは本筋から外れるのでここでは行わない。}。
つまりこれはmod \(N\)の世界の和と積として知られていたものになっている。

これにて mod \(N\) というのが、
2つの整数\(a,b\)が\(N\)で割った余りが同じときに、
これらの元を同一視できるという同値関係を入れることによって得られる\(N\)個の同値類
\([0],\ldots,[N-1]\)へと整数をグループ分けして、
これらのグループでの演算を定義することによって得られる世界\(\mathbb{Z}_N\)であることがわかった。

次に和と密接にかかわる演算である差を定義することを考える。
その手掛かりとして、引き算は任意の元\(a\)に対して、自分自身との差が0つまり
\begin{equation}\begin{aligned}[b]
    a - a = 0
\end{aligned}\end{equation}
という関係が成り立つと我々は考えれるだろう。
そうしたときに、この式で唯一現れる具体的な数である\(0\)に注目してみる。
この\(0\)という元は和に現れたときには、
任意の元\(a\)に対して
\begin{equation}\begin{aligned}[b]
    a+0 = 0+a = a
\end{aligned}\end{equation}
というように元の値を変えないような性質を持っている。
このような性質を持つ元のことを和に関する単位元と呼ぶ。
つまり\(\mathbb{Z}_N\)において\([0]\)は和に関する単位元である。

改めて、\(a-a=0\)の式に戻る。この式をあえて和の式として書き換えると
\begin{equation}\begin{aligned}[b]
    a + (-a) = 0
\end{aligned}\end{equation}
というように書くことができる。
こうして書くと\((-a)\)というのは新たな元とみなせて、
この元の性質として\(a\)に足すことで0になるようなものであるとわかる。
この\(a\)に対する\((-a)\)のことを\(a\)の逆元と呼ぶ。
つまり引き算\(a-b\)は\(a\)に\(b\)の逆元を足す演算のことである。
\(\mathbb{Z}_N\)において、\([a]\)の逆元は
\begin{equation}\begin{aligned}[b]
    [a]+[N-a]=[a+(N-a)]=[N]=[0]
\end{aligned}\end{equation}
であることから、\([N-a]\)であるとわかる。
また、代表元の自由度を利用して、\([N-a]\)の代表元を\(N-a\)ではなく\(-a\)と書いてもよいことがわかる。
逆元がわかったので、差は
\begin{equation}\begin{aligned}[b]
    [a]-[b] := [a]+(-[b]) = [a]+[N-b] = [a+(N-b)] = [a-b]
\end{aligned}\end{equation}
のように定義されることがわかる。

このように集合に2項演算が入っていて単位元や逆元が存在するものは多くある。
これに加えて2項演算に更に良い性質が少し入ったものには群と名付けられている。
以下に数学的な定義を書いてみよう。
\begin{NoteBox}[title=群の定義]{}
    ある集合\(S\)に対して、二項演算\(+\)が定義されているとする。
    このとき、任意の\(a,b,c\in S\)に対して以下の条件を満たすとき、\(S\)は群であると言う。
    \begin{itemize}
        \item 結合律: \(a+(b+c) = (a+b)+c\).
        \item 単位元の存在: \(S\)の中に\(0\)という元があって、任意の\(a\in S\)に対して\(a+0=0+a=a\)となる。
        \item 逆元の存在: 任意の\(a\in S\)に対して、\(S\)の中に\(-a\)という元があって、\(a+(-a)=(-a)+a=0\)となる。
    \end{itemize}
    また、さらに任意の\(a,b \in S\)に対して\(a+b=b+a\)が成り立つとき、\(S\)は可換群であると言う。
\end{NoteBox}
これにより\(\mathbb{Z}_N\)は加法に関して可換群であると言える。

それでは積の方についても同様に単位元や逆元を考えてみよう。
積に関する単位元は、任意の元\(a\)に対して
\begin{equation}\begin{aligned}[b]
    a\cdot 1 = 1 \cdot a = a
\end{aligned}\end{equation}
という性質を持つ\(1\)という元である。
\(\mathbb{Z}_N\)において、\([1]\)は積に関する単位元である。

積に関する逆元は、和の単位元\(0\)を除く任意の元\(a\)に対して
\begin{equation}\begin{aligned}[b]
    a \cdot a^{-1} = a^{-1} \cdot a = 1
\end{aligned}\end{equation}
という性質を持つ\(a^{-1}\)という元である。
\(0\)の逆元は整数\(\mathbb{Z}\)や実数\(\mathbb{R}\)であるように、
整合的な定義は存在しないので除外したということである。

ではこの積に関する逆元は\(\mathbb{Z}_N\)の中に存在するだろうか。
すこし例を考えてみる。
\begin{table}
    \centering
    \caption{\(\mathbb{Z}_N\)の積に関する演算表}
    \label{table:chap2_ZN_product}
    \begin{minipage}{0.4 \columnwidth}
        \centering
        \subcaption{\(\mathbb{Z}_5\)}
        \label{table:chap2_Z5_product}
        \begin{tabular}{c|ccccc}
            \hline
            \(\cdot\) & [0] & [1] & [2] & [3] & [4]\\
            \hline
            [0] & [0] & [0] & [0] & [0] & [0]\\ \relax
            [1] & [0] & [1] & [2] & [3] & [4]\\ \relax
            [2] & [0] & [2] & [4] & [1] & [3]\\ \relax
            [3] & [0] & [3] & [1] & [4] & [2]\\ \relax
            [4] & [0] & [4] & [3] & [2] & [1]\\
            \hline
        \end{tabular}
    \end{minipage}
    \begin{minipage}{0.4 \columnwidth}
        \centering
        \subcaption{\(\mathbb{Z}_6\)}
        \label{table:chap2_Z6_product}
        \begin{tabular}{c|cccccc}
            \hline
            \(\cdot\) & [0] & [1] & [2] & [3] & [4] & [5]\\
            \hline
            [0] & [0] & [0] & [0] & [0] & [0] & [0]\\ \relax
            [1] & [0] & [1] & [2] & [3] & [4] & [5]\\ \relax
            [2] & [0] & [2] & [4] & [0] & [2] & [4]\\ \relax
            [3] & [0] & [3] & [0] & [3] & [0] & [3]\\ \relax
            [4] & [0] & [4] & [2] & [0] & [4] & [2]\\ \relax
            [5] & [0] & [5] & [4] & [3] & [2] & [1]\\
            \hline
        \end{tabular}
    \end{minipage}
\end{table}
表\ref{table:chap2_ZN_product}はそれぞれ\(\mathbb{Z}_5\)と\(\mathbb{Z}_6\)における積に関する演算表である。
表\ref{table:chap2_Z5_product}の演算表を見ると、
\([1]\)の逆元は\([1]\),
\([2]\)の逆元は\([3]\),
\([3]\)の逆元は\([2]\),
\([4]\)の逆元は\([4]\)であることがわかる。
つまりちゃんと逆元が存在していることがわかる。
一方、表\ref{table:chap2_Z6_product}の演算表を見ると、
\([1]\)の逆元は\([1]\),
\([5]\)の逆元は\([5]\)であることがわかる。
しかし、\([2],[3],[4]\)の逆元は存在しないことがわかる。

このように、和については群をなし、
積についても単位元が存在して結合律や分配律といった基本的な計算規則は満たしているものの、
積の逆元が必ずしも存在しない集合のことを環（Ring）と呼び、
\(0\)以外の元全てに対して積の逆元が存在する集合のことを体（Field）と呼ぶ。
数学の形式的な定義は以下の通りである。
\begin{NoteBox}[title=環と体の定義]{}
    ある集合\(S\)に対して、二項演算\(+\)と\(\cdot\)が定義されているとする。
    このとき、任意の\(a,b,c\in S\)に対して以下の条件を満たすとき、\(S\)は環であると言う。
    \begin{itemize}
        \item \(S\)は二項演算\(+\)について可換群である。
        \item 積の結合律: \(a\cdot (b\cdot c) = (a\cdot b)\cdot c\).
        \item 積の単位元の存在: \(S\)の中に\(1\)という元があって、任意の\(a\in S\)に対して\(a\cdot 1=1 \cdot a = a\)となる。
        \item 分配律: \(a \cdot (b+c) = a \cdot b + a \cdot c,\, (a+b)\cdot c = a \cdot c + b \cdot c\).
    \end{itemize}
    また、さらに任意の\(a,b \in S\)に対して\(a\cdot b = b \cdot a\)が成り立つとき、\(S\)は可換環であると言う。
    さらに、\(0\)以外の任意の元\(a\in S\)に対して、\(a^{-1}\)という元があって、
    \begin{equation}\begin{aligned}[b]
        a \cdot a^{-1} = a^{-1} \cdot a = 1
    \end{aligned}\end{equation}
    となるとき、\(S\)は体であると言う。
\end{NoteBox}
今まで出てきた、群・環・体を平たく言うと、
足し算・引き算ができるのが群、
足し算・引き算・掛け算ができるのが環、
足し算・引き算・掛け算・割り算ができるのが体であると言える。

\(\mathbb{Z}_N\)は積に関しては単位元が存在して結合律や分配律といった基本的な計算規則は満たしているものの、
積の逆元が必ずしも存在しないため、可換環であると言える。
さらに\(N\)が特定の場合には0以外の全ての元に対して積の逆元が存在するため、体であると言える。
\(\mathbb{Z}_N\)が体であるための条件は、\(N\)が素数\(p\)であることである
\footnote{これはフェルマーの小定理（
\(a^{p-1} \equiv 1 \pmod p\) より \(a \cdot a^{p-2} \equiv 1\) となるため、
\(a^{p-2}\) が逆元となる）などからわかる。}。
そうした体のことをガロア体と呼び、\(\mathbb{F}_{p}\)のように書かれることが多い。

また、物理における\(\mathbb{Z}_N\)の慣例を紹介する。
統計力学や場の量子論などで「\(\mathbb{Z}_N\) 対称性」といった言葉を耳にしたことがあるだろう。
例えばイジングモデルが持つすべてのスピンをひっくり返しても同じエネルギーとなるという対称性は \(Z_2\) と呼ばれる。
物理の文脈で \(\mathbb{Z}_N\) と書かれるとき、
多くの場合、この「和に関する群」としての性質のみに注目している。
ある操作を \(N\) 回繰り返すと元に戻るような周期的な構造を持つ群を、
巡回群（Cyclic group）と呼び、物理で現れる \(\mathbb{Z}_N\) はまさにこの巡回群のことを指している。
このとき、積に関しては特に定義しないことが多い。
符号理論における\(\mathbb{Z}_N\)は環のことを指しているので、
物理とnotationが衝突していることに注意が必要である。
そのため、以降では符号理論や環の方を黒板太字の\(\mathbb{Z}_N\),
物理の巡回群の方を通常の\(Z_N\)と書いて区別することにする。

ここでこのセクションのはじめで触れた mod 2 の特異性について振り返ってみよう。
2は素数である。したがって、mod 2の世界は単なる環ではなく、
四則演算が完備された立派な体 \(\mathbb{F}_2\) である。
\(\mathbb{F}_2\) は元が \(0\) と \(1\) の2つしか存在しないため、
加法の単位元、加法の逆元、積の単位元といった代数構造の根幹を、
この2つの元だけで全て担わなければならない。
その結果、\(1+1=0\) すなわち \(1\) の加法の逆元が \(1\) 自身となるという一見奇妙な性質が現れていたのである。
以降ではこれを避けるため\(\mathbb{F}_p\)での計算を用いて、
量子誤り訂正符号を記述していくことになる。


\section{一般化パウリ演算子とqudit}
bit の世界である \(\mathbb{F}_2\) における量子ビットは \(\ket{0},\ket{1}\)であった。
これから\(\mathbb{F}_p\)を考えるので、これに対応する情報を蓄える状態を考える。
\(\mathbb{F}_p\)はp種類の数があるので、これらの数を表すためには、
\(\ket{0},\ket{1},\ldots,\ket{p-1}\)の\(p\)個の状態を必要とする。
このとき二進数を表す binary からくる bit という名称では不適当であるので、
このような量子ビットを一般化したものを qudit と呼ぶ。

量子状態の次元が変わったので、量子ゲートも変える必要がある。
bit の世界では、Xゲートは \(\ket{0}\) と \(\ket{1}\) を入れ替えるゲートであった。
すなわち状態を別の状態に遷移させる演算子であった。
したがって、\(\mathbb{F}_p\)の世界における一般化されたXゲートも状態を別の状態へと写すゲートとして定義できる。
これをパウリ\(X\)を拡張したものとして
\begin{equation}\begin{aligned}[b]
    \mathcal{X}\ket{j} &= \ket{j+1 \pmod p} \qquad (j=0,1,\ldots,p-1)\\
    \mathcal{X}^\dagger \ket{j} &= \ket{j-1 \pmod p} \qquad (j=0,1,\ldots,p-1)
\end{aligned}\end{equation}
というものを考える。これを一般化パウリ\(X\) やシフト演算子と呼ぶ。
\(\ket{0},\ket{1},\ldots, \ket{p-1}\)をそれぞれ
列ベクトル\((1,0,\ldots, 0)^\mathsf{T},\) \((0,1,\ldots, 0)^\mathsf{T},\) \((0,0,\ldots, 1)^\mathsf{T},\)
のように表すとすると、この一般化パウリ\(X\)およびそのエルミート共役の行列表示は
\begin{equation}\begin{aligned}[b]
    \mathcal{X}_p &= \begin{pmatrix}
        0 & 0 & \cdots & 0 & 1 \\
        1 & 0 & \cdots & 0 & 0 \\
        0 & 1 & \cdots & 0 & 0 \\
        \vdots & \vdots & \ddots & \vdots & \vdots \\
        0 & 0 & \cdots & 1 & 0
    \end{pmatrix}, &
    \mathcal{X}_p^\dagger = \begin{pmatrix}
        0 & 1 & 0 & \cdots & 0 \\
        0 & 0 & 1 & \cdots & 0 \\
        \vdots & \vdots & \vdots & \ddots & \vdots \\
        0 & 0 & 0 & \cdots & 1\\
        1 & 0 & 0 & \cdots & 0
    \end{pmatrix}
\end{aligned}\end{equation}
となる。\(p=2\)のときはよく知っている\(X\)の行列表示となり、確かに拡張したものであるとわかる。
さらに\(p\geq 3\)の場合は\(\mathcal{X}_p\neq \mathcal{X}_p^\dagger\)であるのがわかる。
これは状態のラベルを増やす操作と状態のラベルを減らす操作を明確に区別することにつながっている。

同様にパウリ\(Z\)の一般化を行おう。
bit の世界では、Zゲートは \(\ket{0}\)には位相\(+1\)を、
\(\ket{1}\)には位相\(-1\)をつけるゲートであり、
すなわち現在の状態のインデックスに応じて異なる位相をつけるゲートであった。
この性質から\(\mathbb{F}_p\)における一般化されたZゲートも状態のインデックスに応じて異なる位相をつけるゲートとして定義できる。
\begin{equation}\begin{aligned}[b]
    \mathcal{Z}_p\ket{j} &= \omega_p^j \ket{j} \qquad (j=0,1,\ldots,p-1)\\
    \mathcal{Z}_p^\dagger \ket{j} &= \omega_p^{-j} \ket{j} \qquad (j=0,1,\ldots,p-1)\\
    \omega_p &= e^{2\pi i/p}
\end{aligned}\end{equation}
というものを考える。これを一般化パウリ\(Z\) やクロック演算子と呼ぶ。
これの行列表示は対角行列となり、
\begin{equation}\begin{aligned}[b]
    \mathcal{Z}_p &= \begin{pmatrix}
        1 & 0 & 0 & \cdots & 0 \\
        0 & \omega_p & 0 & \cdots & 0 \\
        0 & 0 & \omega_p^2 & \cdots & 0 \\
        \vdots & \vdots & \vdots & \ddots & \vdots \\
        0 & 0 & 0 & \cdots & \omega_p^{p-1}
    \end{pmatrix}, &
    \mathcal{Z}_p^\dagger &= \begin{pmatrix}
        1 & 0 & 0 & \cdots & 0 \\
        0 & \omega_p^{-1} & 0 & \cdots & 0 \\
        0 & 0 & \omega_p^{-2} & \cdots & 0 \\
        \vdots & \vdots & \vdots & \ddots & \vdots \\
        0 & 0 & 0 & \cdots & \omega_p^{-(p-1)}
    \end{pmatrix}
\end{aligned}\end{equation}
と表される。\(p=2\)のときは\(\omega_2=-1\)なので、
\(\mathcal{Z}_p\)はよく知っている\(Z\)の行列表示と一致する。

このように定義された\(\mathcal{X}_p,\mathcal{Z}_p\)は通常のパウリ\(X,Z\)と同様に、
互いに交換しない性質を持っている。
通常のパウリとは異なり\(\mathcal{X}_p^a,\mathcal{Z}_p^b\)のような冪も考えることができる。
この交換関係を調べると
\begin{equation}\begin{aligned}[b]
    \mathcal{Z}_p^b \mathcal{X}_p^a \ket{j} &= \omega_p^{bj} \ket{j+a },\\
    \mathcal{X}_p^a \mathcal{Z}_p^b \ket{j} &= \omega_p^{b(j+a)} \ket{j+a}\\
    \Rightarrow \mathcal{X}_p^a \mathcal{Z}_p^b &= \omega_p^{ab} \mathcal{Z}_p^b \mathcal{X}_p^a
\end{aligned}\end{equation}
のようにわかる。ここに\(a=b=1, p=2\)として通常のパウリと同じ結果になるかを確認すると
\begin{equation}\begin{aligned}[b]
    \mathcal{X}_2 \mathcal{Z}_2 &= \omega_2^{1\cdot 1} \mathcal{Z}_2 \mathcal{X}_2 = -1 \cdot \mathcal{Z}_2 \mathcal{X}_2\\
    \Rightarrow 0 &= \mathcal{X}_2 \mathcal{Z}_2 + \mathcal{Z}_2 \mathcal{X}_2 = \{\mathcal{X}_2,\mathcal{Z}_2\}
\end{aligned}\end{equation}
となり確かに通常のパウリの持つ反交換関係\(\{X,Z\} = 0\)が得られる。

\section{Shor符号の一般化}


\end{document}
